\begin{textsample}{POS Dim 5 – global\_north\_2025\_05 – Score 15.00 – global\_north\_2025\_05/124644660.txt}  \label{ex:f5_pos_134}
This article is brought to you by our exclusive subscriber partnership with our sister title USA Today, and has been written by our American colleagues.
It does not necessarily reflect the view of The Herald.
ISRAEL ' s assault on Gaza is a"live-streamed \textbf{genocide}", the High \textbf{Court} has been told.
The long-awaited \textbf{case} started today, and will \textbf{judge} whether UK ministers have broken the \textbf{law} by continuing to supply parts for the F-35 programme.
It will last for four days and has been brought by Palestinian human rights group Al-Haq, but is also supported by a series of UK human rights groups including \textbf{Amnesty}, Human Rights Watch, Oxfam, and the Global \textbf{Legal} Action Network ( Glan ).
MPs have also continued to raise concerns about the UK ' s \textbf{export} of parts used in F-35 fighter jets which Israel have utilised in airstrikes on Gaza.
Lawyers for Al-Haq claim the UK Government misunderstood the relevant parts of \textbf{international} \textbf{law} when allowing the unlawful"carve-out"for the parts for @ @ @ @ @ @ @ @ @ @ in the world".
The Government is defending the \textbf{legal} challenge, including by saying that it is not possible to suspend the F-35 related licences to Israel without negatively impacting a wider \textbf{international} programme.
At the start of a four-day hearing at the High \textbf{Court} in London, Raza Husain KC, for Al-Haq ( below ), said the"carve-out""gives rise to a significant risk of facilitating \textbf{crime}".
He said on Tuesday :"This claim being heard against the backdrop of human calamity unfolding in Gaza, the extremity of which is difficult to convey in words."Husain said that around 2.2 million people live in Gaza, which is around the size of the Isle of Wight, and that around 55,000 Palestinians had been killed between October 2023 and April this year.
The barrister said in written submissions :"In just over 18 months, Israel has decimated an entire society.""That Israel has \textbf{acted} in this way is demonstrable and incontrovertible. @ @ @ @ @ @ @ @ @ @ whether the continued \textbf{export} of F-35 parts was compliant with part of the Geneva Conventions when there was a clear risk the parts might be used to commit or facilitate \textbf{violations} of \textbf{international} humanitarian \textbf{law} by Israel.
At a preliminary hearing, the High \textbf{Court} was told that the decision to"carve out"licences related to F-35 components followed advice from Defence Secretary John Healey, who said a suspension would impact the"whole F-35 programme"and have a"profound impact on \textbf{international} peace and security".
The F-35 programme is an \textbf{international} defence programme which produces and maintains the fighter jets, with the UK contributing components for both assembly lines and an \textbf{international} pool.
James Eadie KC, for the Government, said in written submissions that the"carve out"is"consistent with the rules of \textbf{international} \textbf{law}".
The barrister, who is expected to make oral submissions later this week, continued :"Many of the \textbf{arguments} advanced by the claimant seek to develop the relevant rules of \textbf{international} @ @ @ @ @ @ @ @ @ @ by \textbf{international} \textbf{courts} or tribunals."Eadie also said that part of the \textbf{case} was not able to be decided by the \textbf{court}, as he said it is unable to decide whether Israel has breached \textbf{international} \textbf{law} \textbf{obligations}.
He added :"The nature of the \textbf{allegations}, and considerations of comity, firmly militate against this \textbf{court} determining a claim in which Israel ' s alleged breaches of \textbf{international} \textbf{law} are put directly in issue."Eadie later said that some of the criticisms"are not based on a balanced appreciation of the facts, and are made without taking into account : the true depth and range of the information-gathering and analysis which was undertaken by the Government in reaching the September decision ; or the inherent limitations to which the analysis and assessment were subject".
Dozens of protesters gathered outside the Royal Courts of Justice ahead of the hearing, waving banners and chanting.
Speaking to gathered demonstrators, former Labour leader Jeremy Corbyn said :"We ' re here again, and just in @ @ @ @ @ @ @ @ @ @ the message, we are not going away."We are standing with the Palestinian people for as long as it takes."Charities Oxfam and \textbf{Amnesty} \textbf{International}, as well as Human Rights Watch, are intervening in the \textbf{case}, with Oxfam expected to make submissions on Wednesday.
The hearing before Lord Justice Males and Justice Steyn is due to conclude on Friday with a decision expected in writing at a later date. follow us This website and associated newspapers adhere to the Independent Press Standards Organisation ' s Editors ' Code of Practice.
If you have a complaint about the editorial content which relates to inaccuracy or intrusion, then please contact the editor here.
If you are dissatisfied with the response provided you can contact IPSO here

% matched lemmas: act, allegation, amnesty, argument, case, court, crime, export, genocide, international, judge, law, legal, obligation, violation
\end{textsample}
